% Synthetic LaTeX for the citation-extraction tests (spec section 17). CC0-1.0.
%
% This document exists to be a paper that *refers to other papers*, which none of the other
% fixtures do. It carries one of each kind of reference the extractor is allowed to match —
% identifier, author-and-year, near-complete title — and several that it must refuse.
%
% Unlike the other fixtures, its manifest entry is keyed to a paper that actually exists in
% `fixtures/papers.sample.json` (`arxiv:2604.14247`), and every paper it refers to is in that
% corpus too. That is what lets `services/mentions` produce real relations from `make seed`
% rather than from a document whose references point at nothing.
\documentclass{article}
\usepackage{amsmath}

\title{Large Deviations in Log-Concave Models}
\author{M. Zubkov \and F. Castellanos \and R. Xu \and K. Sasaki}

\begin{document}
\maketitle

\section{Introduction}

% Refers by arXiv identifier: the strongest form the extractor recognises.
% Target: "Concentration for Non-Reversible Measures" (2026), arXiv:2612.10959.
The evaluation in arXiv:2612.10959 obtains the deviation constant by a spectral argument,
and we follow that route here.

% Refers by surname and year together, and draws a contrast. Neither half of the reference
% may match alone, and the contrastive cue is what lets `relations.classify` reach the
% `contrasting` label at all.
% Target: "Thresholds for Sparse Subgraphs" (2026), R. Almeida.
In contrast to Almeida (2026), we do not require the threshold to be sharp, and the bound
therefore extends to the degenerate case.

% Refers by a nearly complete title, and the title is long enough to be distinctive.
% Target: "Improved Bounds for Off-Diagonal Ramsey Numbers" (2025).
The paper Improved Bounds for Off-Diagonal Ramsey Numbers treats the multi-dimensional case.

% A title too short to be distinctive. Must not match: "Sublinear Geometric Estimation"
% contributes three usable words, and any sentence about estimating geometry contains them.
The paper Sublinear Geometric Estimation treats only the one-dimensional case.

\section{Discussion}

% A year with no surname. Must not match anything.
Earlier treatments from 2026 used a different normalisation.

% A surname with no year. Must not match anything.
The argument of Almeida is standard.

% A contrastive phrase in a sentence that refers to nobody. Must not attach to any paper.
Unlike the usual presentation, we keep the constant explicit throughout.

\end{document}
